% sections/05_comments.tex  (TV7)

\section{Comments on the topic}\label{sec:comments}\writer{TV7}

\subsection{What has been accomplished}

We compared our work with the requirements of the project \cite{zemlyanova2020}:
\begin{itemize}
  % CHECK BEFORE SUBMISSION (TV1's review, point 3): the sentence "We finished all
  % 14 tasks ... answered Q1--Q6" is only true when the group has really finished
  % all 14 tasks. At the time of the review only 8/14 were done. Do not change the
  % wording now; re-check it against the final work right before submitting.
  \item \textbf{Code:} We finished all 14 tasks in one Python file,
    \texttt{lab03.py}, and kept the same variable names as in the project
    (\texttt{ImJPG}, \texttt{ImJPG\_User}, \texttt{ImJPG\_UserInv}, \dots).
    All matrix filters use the same loop over pixels as the project code
    (Section~\ref{sec:code}).
  \item \textbf{Questions:} We answered Q1--Q6 in the code and explained them
    again in Section~\ref{sec:ex-material}.
  \item \textbf{Checking:} Our results match the material: the image size is
    $308 \times 322 \times 3$, $\det(\texttt{UserMatrix}) = 0.3025$, and the
    inverse filter recovers \texttt{ImJPG\_User} with a maximum error of 1.
  \item \textbf{Theory:} The report explains why a filter is a matrix
    multiplication, when it can be inverted, which filters lose information, and
    why the sepia filter cannot be inverted in practice.
\end{itemize}

\subsection{Limitations}

\begin{itemize}
  \item We only studied filters that change each pixel by itself. Filters that use
    the pixels around it, like blurring or sharpening, are not included.
  \item An image only stores whole numbers from 0 to 255. Rounding and clipping
    make the inverse only approximate, and we studied this mostly by experiments,
    not by a general formula.
  \item Filtering with a loop over every pixel, as in the project, is slow in
    Python (about one second for each filter, even for a small image).
  \item We used Python instead of MATLAB. Because of rounding, a few pixels may be
    different from the MATLAB result by 1 level.
  \item We did not try any better method to recover the sepia image.
  \item The results were checked on only a small number of images, and we judged
    the quality mostly by eye.
\end{itemize}

\subsection{Possible further development}

\begin{itemize}
  \item \textbf{Affine filters:} Brightness changes and color inversion (Task~9)
    add a constant, so they are not linear. Writing a color as $[R;G;B;1]$ and
    using a $4 \times 4$ matrix puts them into one matrix multiplication.
  \item \textbf{Better recovery:} For an almost singular filter like sepia, we
    could use least squares or regularization instead of the direct inverse
    \cite{higham2002}.
  \item \textbf{Neighbor filters:} Convolution with a small kernel can blur or
    sharpen an image or find edges \cite{gonzalez2018}.
  \item \textbf{Other color spaces:} In YCbCr or HSV, brightness and color are
    separated, so changing the color without changing the brightness is easier
    \cite{gonzalez2018}.
  \item \textbf{A small application:} The user types a filter matrix and sees the
    new image together with its determinant and condition number.
\end{itemize}
